\section{Operator positivity through adjoint sandwiches}
\label{sec:operator}
\status{Independent operator replay, an automatic homogeneous Gram producer,
and a bounded finite-ray producer are implemented. General SDP discovery and
Lean positivity reconstruction are integration work, not executed features.}

\subsection{The correct replacement for scalar constraint products}

For real symmetric positive semidefinite $G$ and any real matrix $Q$,
\[
 z^{\mathsf T}(Q^{\mathsf T}GQ)z=(Qz)^{\mathsf T}G(Qz)\ge0.
\]
Thus congruence preserves positivity. In contrast, let
\[
 A=\begin{pmatrix}1&0\\0&0\end{pmatrix},\qquad
 B=\begin{pmatrix}1&1\\1&1\end{pmatrix}.
\]
Both are positive semidefinite, but $AB$ is not symmetric. Even the symmetric
product $AB+BA$ is not positive semidefinite: at $z=(1,-2)^{\mathsf T}$ its
quadratic form is $-2$. This exact arithmetic example is an executed regression
test, not just a caution in the prose.

The operator receipt therefore has the form
\begin{equation}\label{eq:operator-cert}
 p=\sum_{k=1}^{t} d_k q_k^*g_{i_k}q_k
    +\sum_{\ell=1}^{u}c_\ell a_\ell r_{j_\ell}b_\ell,
 \qquad d_k\in\Q_{\ge0},\quad g_0=1.
\end{equation}
The original problem supplies $g_i\succeq0$ and $r_j=0$. The equality terms have
no sign restrictions and preserve their left/right word contexts. Here the
$q_k$ are polynomials, whereas $a_\ell,b_\ell$ are words; polynomial equality
multipliers can be expanded into finitely many such word contexts.

\begin{theorem}[Operator receipt soundness]\label{thm:operator}
A checked identity~\eqref{eq:operator-cert} implies $p(X)\succeq0$ for every
matrix tuple satisfying the original involution, equality, and positivity
assumptions.
\end{theorem}
\begin{proof}
The equality contribution evaluates to zero by the same argument as
Theorem~\ref{thm:ideal-sound}. Each remaining summand is positive semidefinite
by congruence and nonnegative scalar multiplication. Their sum is positive
semidefinite. The involution contract makes the formal $q_k^*$ evaluate to the
transpose of the evaluated $q_k$.
\end{proof}

The proof does not assert that this certificate language expresses every valid
conditional matrix inequality. In particular, a finite dictionary of $q_k$ is a
further restriction. Failure of that search is not a mathematical counterexample.

\subsection{A projection example that requires equality cooperation}

Suppose $x^*=x$ and $x^2=x$. Then
\[
 x=x^*x-(x^2-x)
\]
is a one-square operator receipt. The target has degree one, so it is not in the
free homogeneous-square fragment below. The equalities are what make it
positive. The implemented finite-ray worker finds the weight and the ideal
residual rather than trying to treat the premise $x^2-x=0$ as a positivity
constraint.

Similarly, from $x\succeq0$ the target $yxy$ with self-adjoint $y$ is certified
by the sandwich $y^*xy$. Deleting the positive premise makes the exact same
receipt invalid. Both cases are covered by the stored worked examples and
boundary tests.

\section{Automatic homogeneous squares at arbitrary even degree}
\label{sec:gram}

\subsection{Why the noncommutative Gram matrix is unique}

Fix $d\ge0$ and list all words of length $d$ as
$W_d=(w_1,\ldots,w_N)^{\mathsf T}$, where $N=m^d$. A homogeneous polynomial
$p$ of degree $2d$ has a unique expansion
\begin{equation}\label{eq:gram}
 p=W_d^* Q W_d=\sum_{i,j}Q_{ij}w_i^*w_j.
\end{equation}
Indeed, every word of length $2d$ has a unique split into its first and last
$d$ letters; the first half uniquely determines $w_i$ by reversing and applying
the atom involution. There are $m^{2d}=N^2$ words and $N^2$ ordered matrix entries.
Formal self-adjointness of $p$ is exactly symmetry of the rational matrix $Q$.

This is a decisive difference from ordinary commuting multivariate polynomials,
where different products of monomials can have the same exponent vector and
Gram representations need not be unique. The existing scalar quadratic $LDL$
mechanism does not, by itself, expose this arbitrary-even-degree free-word
fragment. The new step is the word splitting and its exact coefficient
correspondence, not the invention of $LDL$ factorization.

\begin{theorem}[Complete homogeneous-square fragment]\label{thm:gram}
For a rational homogeneous formally self-adjoint $p$ of degree $2d$, the
following are equivalent:
\begin{enumerate}[label=(\roman*),nosep]
\item Its unique word Gram matrix $Q$ is positive semidefinite.
\item $p=\sum_k d_k q_k^*q_k$, with rational $d_k\ge0$ and rational
homogeneous $q_k$ of degree $d$.
\item $p$ has a sum-of-Hermitian-squares representation using real homogeneous
polynomials of degree $d$.
\end{enumerate}
The rational representation can be found by exact pivoted $LDL^{\mathsf T}$
elimination, without numerical semidefinite optimization.
\end{theorem}
\begin{proof}
A representation in (iii) expresses $Q$ as a sum of real outer products by
uniqueness in~\eqref{eq:gram}, proving (i). Condition (ii) immediately implies
(iii), absorbing $\sqrt{d_k}$ into the real polynomials.

For (i) to (ii), work entirely over $\Q$. If a diagonal entry $Q_{kk}=a$ is
positive, complete the square using the row
\[
 \ell=e_k+\sum_{j\ne k}\frac{Q_{kj}}{a}e_j.
\]
Subtract $a\ell\ell^{\mathsf T}$ and recurse on the remaining Schur complement.
It is positive semidefinite and rational. If all remaining diagonal entries are
zero, positive semidefiniteness forces every off-diagonal entry to be zero: a
nonzero $Q_{ij}$ would give a negative quadratic form on a suitable vector
supported on $i,j$. Thus the process terminates with
$Q=\sum_k d_k\ell_k\ell_k^{\mathsf T}$, $d_k\in\Q_{>0}$. Set
$q_k=\sum_i(\ell_k)_i w_i$ and use~\eqref{eq:gram}.
\end{proof}

The zero polynomial is accepted with an empty square list. A nonzero polynomial
of odd homogeneous degree is outside this entrance fragment. The matrix test is
exact: an off-diagonal nonzero entry with zero diagonal is a rejection, not an
invitation to add a numerical tolerance.

\subsection{Relation to matrix positivity, without overstating the software}

Helton's theorem identifies matrix-positive free polynomials in self-adjoint
variables with sums of Hermitian squares~\cite{helton}. It supplies relevant
mathematical context; the present report does not claim that theorem as new.
For a homogeneous target, square factors can be taken homogeneous of half the
degree. To see the degree restriction, consider the highest homogeneous pieces
of putative factors: their sum of squares cannot cancel, because the unique
word Gram matrix is a sum of outer products and is nonzero if any piece is
nonzero. Apply the same argument to the lowest nonzero pieces. A pure
 degree-$2d$ sum therefore has no factor degrees above or below $d$.

Consequently, for self-adjoint free variables and all finite matrix dimensions,
the exact Gram criterion also characterizes homogeneous matrix positivity.
That mathematical corollary does \emph{not} mean the delivered negative output is
a Lean proof of nonpositivity. The routine returns no certificate when $Q$ is
not positive semidefinite. It does not construct a self-adjoint matrix tuple
violating the inequality, and the checker never accepts a Gram failure as such
a countermodel. The separately implemented quotient-matrix construction is an
\emph{equality} countermodel mechanism with a different scope.

\subsection{An automatic degree-six example}

Let self-adjoint $x,y$ be free and set
\[
 q=x^2y+2yx^2-3y^3,\qquad p=q^*q.
\]
The experiment supplies only the expanded target $p$ to
\code{homogeneous\_sos}. There is no supplied factor or dictionary. The
producer builds an $8\times8$ Gram matrix, factors it rationally, and emits a
square receipt which independent replay expands back to $p$. The recovered
factor need not have the same normalization as the generating $q$; the
nonnegative rational weight accounts for that harmless freedom.

The cost is $O(m^{3d})$ dense rational field operations in the worst case, with
an $m^{2d}$-entry Gram matrix. Sparse word support and block decomposition may
reduce that cost; a six-variable degree-twelve goal is not made small merely by
being in a decidable fragment. A production implementation needs explicit word,
coefficient-bit, and factor-size budgets.

\section{Mixed-degree search modulo equalities}
\label{sec:finite-rays}

\subsection{Quotient the equality span before solving for positive weights}

Fix a context bound $D$, form the equality column span $M_D$ from
Section~\ref{sec:ideal}, and choose a finite family of candidate rays
\[
 a_k=q_k^*g_{i_k}q_k.
\]
Let $\pi_D$ be the linear remainder map of the exact echelon basis. Search for
\begin{equation}\label{eq:ray}
 \pi_D(p)=\sum_k d_k\pi_D(a_k),\qquad d_k\ge0.
\end{equation}
After obtaining the weights, reduce the \emph{original} residual
$p-\sum_kd_ka_k$ to recover its two-sided ideal receipt. This avoids presenting
all equality coefficients as free signed variables in the positive feasibility
problem.

\begin{proposition}[Equality projection loses no fixed-language solution]
For the fixed equality span and fixed rays, equation~\eqref{eq:ray} is feasible
if and only if there exists a receipt~\eqref{eq:operator-cert} using those rays
and equality contexts in $M_D$.
\end{proposition}
\begin{proof}
The kernel of the linear remainder map is precisely $M_D$. Thus the projected
identity holds exactly when the difference of target and weighted rays belongs
to $M_D$, which is exactly what the equality terms express.
\end{proof}

It would be incorrect to run positivity directly on the possibly
non-self-adjoint-looking remainder and forget the target. A normal-form map is
only a linear coordinate device; the final replay always checks the original
self-adjoint target, the original positive premises, and the original relations.

\subsection{What the implemented support search guarantees}

The standard-library producer enumerates supports of size one through a given
cap, solves the corresponding exact linear systems, and accepts a representation
only if every recovered coefficient is nonnegative. It defaults to support
size at most four and at most 10,000 attempted supports. It is not an SDP solver.

There is a precise completeness statement for an uncapped finite-ray version.
If a vector lies in the cone generated by finitely many vectors in a rational
space of dimension $r$, it has a representation on a linearly independent
support of size at most $r$. Starting with a nonnegative representation, a
linear dependence of its positive-support vectors allows movement along a
coefficient null vector until a positive coefficient first becomes zero,
without making any coefficient negative. Repeat until the support is independent.
An independent rational support has unique rational coefficients for a rational
target. Hence exhaustive support enumeration through $r$ finds a solution when
one exists. With a smaller support or attempt cap, failure remains \unknown{}.

Dependent supports are not a loophole in this argument: their feasible
representations can be reduced to an independent support, which is eventually
visited by the uncapped enumerator. The result is completeness for this
\emph{fixed cone}, not for conditional noncommutative positivity.

\subsection{Automatic proposals from involution and commutation relations}
\label{sec:graph-dictionary}

The delivered function \code{involution\_dictionary} adds a concrete source of
rays rather than leaving every $q_k$ to a user. It recognizes nonzero scalar
multiples of
\[
 x_i^2-1,\qquad x_ix_j-x_jx_i
\]
in the supplied relation list. These identify involutive atoms and explicitly
known commuting pairs. For two involutive atoms not explicitly marked commuting,
it proposes both their commutator and their anticommutator. For two disjoint
pairs whose cross-pairs are explicitly commuting, it also proposes the sum and
difference of their commutators.

\begin{lstlisting}
U := atoms with a recognized relation x_i*x_i - 1 = 0
E := pairs with a recognized relation x_i*x_j - x_j*x_i = 0
P := pairs in U not already in E
propose x_i*x_j + x_j*x_i and x_i*x_j - x_j*x_i for each pair in P
for disjoint pairs p,q with every cross-pair in E:
    propose comm(p) + comm(q) and comm(p) - comm(q)
deduplicate nonzero scalar multiples in the SEARCH dictionary
send every resulting square proposal through exact weight and identity replay
\end{lstlisting}

Absence of a commutation relation is not used as a proof of noncommutation. It
only makes a pair eligible for proposals. The dictionary does not learn an
unproved relation, and a missed syntactic pattern only reduces search power.
Indirectly derivable involution or commutation facts would have to be supplied
as previously proved lemmas to this entrance. There is no completeness claim
for the ray generator.

At most $P$ eligible pairs give $2P+2\binom P2$ candidates before filtering and
deduplication. The current pair cap is twenty. The algorithm is small enough to
port directly into a Lean-side proposal worker; its correctness need not be
trusted because every eventual receipt is independently replayed.

\section{A worked operator bound: the CHSH expression}
\label{sec:chsh}

\subsection{The mathematical identity}

Let $A,D,B,C$ be self-adjoint involutions and assume every member of $\{A,D\}$
commutes with every member of $\{B,C\}$. Write
\[
 S=A(B+C)+D(B-C),\qquad P=AD-DA,\quad Q=BC-CB,
\]
\[
 T_A=AD+DA,\qquad T_B=BC+CB.
\]
The cross-commutations imply $S^*=S$ and $PQ=QP$. Expanding with the involution
relations gives
\[
 S^2=4I-PQ,\qquad T_A^2-P^2=4I,\qquad T_B^2-Q^2=4I.
\]
Since $P^*=-P$ and $Q^*=-Q$,
\begin{align}
 \tfrac12(P-Q)^*(P-Q)+\tfrac12T_A^*T_A+\tfrac12T_B^*T_B
 &=\tfrac12(-P^2+2PQ-Q^2)+\tfrac12T_A^2+\tfrac12T_B^2\\
 &=4I+PQ=8I-S^2.\label{eq:chsh-sos}
\end{align}
This proves $S^2\preceq8I$, and hence $\|S\|\le2\sqrt2$ in the usual
finite-dimensional Euclidean operator norm. The norm implication follows by
applying $S^*S\preceq8I$ to each vector and taking the supremum over unit
vectors. The recorded artifact certifies the polynomial inequality and a
separate self-adjointness identity; it contains no Lean norm theorem.

This is a worked application of known operator mathematics, not a new physical
bound. Its value here is that a noncommutative receipt can simultaneously use
operator involution, selectively commuting groups, and a positive-square
identity without pretending all four operators commute.

\subsection{What the producer actually discovers}

The source problem supplies four involution relations and four cross-commutation
relations. From those relations the graph rule of
Section~\ref{sec:graph-dictionary} generates six candidate polynomials. The
finite-ray worker searches their nonnegative combinations modulo degree-four
equality contexts. It finds a three-square receipt with rational weights, then
recovers a two-sided ideal residual. Neither the weights nor the list of three
chosen factors is supplied to that search. This is a bounded structurally guided
search, not a general discovery procedure for quantum inequalities.

The unit suite repeats the same construction under all twenty-four permutations
of the atom names. Those are \emph{metamorphic checks of one family}, not
twenty-four independent mathematical discoveries or a twenty-four-goal tactic
benchmark. They test whether pattern recognition or reconstruction accidentally
depends on the particular ordering $A,D,B,C$.

\subsection{The failed first run and its correction}

The first authoring run used the target $8I-S^2$. As a \emph{free} word polynomial,
that expression is not syntactically self-adjoint: $S$ becomes self-adjoint only
through the cross-commutation hypotheses. The conservative operator producer
correctly refused the input. The initial failing log is retained as
\code{results/initial-failure.log}.

The repair was not to turn off the self-adjointness check. The executed operator
target is
\[
 8I-S^*S,
\]
which is formally self-adjoint. A separate exact equality receipt proves
$S-S^*=0$ using the cross-commutation relations. Substitution then gives the
stated squared bound. This is an example of a useful Forge obligation exchange:
an algebra worker proves the missing structural fact while the positivity
worker retains its simpler, safer input language.

For this particular norm bound there is also a simpler discovery opportunity:
if $U=A(B+C)$ and $V=D(B-C)$, then $U^*U+V^*V=4I$ already follows from the
involution relations, and the parallelogram identity gives
$8I-(U+V)^*(U+V)=(U-V)^*(U-V)$. The graph-based three-square certificate is not
minimal. A source-expression-aware sign-flip proposal could find the one-square
version, but that alternative proposer was not implemented or counted in the
recorded run. Retaining the distinction prevents a successful demonstration
from being mislabeled a proof-size optimum.

\section{Trace positivity and explicit cyclic remainders}
\label{sec:trace}

\subsection{A genuinely different certificate kind}

For square matrices, $\tr(AB-BA)=0$. Consequently trace positivity permits
\begin{equation}\label{eq:trace}
 p=\sum_k d_k q_k^*g_{i_k}q_k
   +\sum_\ell c_\ell a_\ell r_{j_\ell}b_\ell
   +\sum_t e_t(u_tv_t-v_tu_t),\qquad d_k\ge0.
\end{equation}
The final sum need not vanish as an operator. Its coefficients can have either
sign. The checker expands every commutator exactly; it does not trust a flag
saying that a residual is ``cyclically zero.''

\begin{theorem}[Trace receipt soundness]
A checked identity~\eqref{eq:trace}, under the original matrix assumptions,
implies $\tr(p)\ge0$.
\end{theorem}
\begin{proof}
The equality terms vanish, the commutator traces vanish by cyclicity, and the
trace of each positive semidefinite sandwich is nonnegative. For real matrices
the last claim follows directly by summing its nonnegative diagonal quadratic
forms against the standard basis.
\end{proof}

\subsection{A finite, constructive cyclic normal form}

Associate each word with the lexicographically least of its cyclic rotations.
The cyclic normal form of a polynomial collects coefficients at these canonical
words. If this normal form is zero, then
\[
 p=\sum_w c_w\bigl(w-\operatorname{cyc}(w)\bigr),
\]
since the coefficients of each rotation class sum to zero. Choose a split
$w=uv$ whose rotation $vu$ is canonical; its contribution is the explicit
commutator $c_w(uv-vu)$. This proves that cyclic reduction to zero can be
converted into the final receipt language without relying on a normalizer's
unexplained verdict.

The implemented search in trace mode uses this linear cyclic map in place of
the equality remainder map. It currently accepts \emph{no equality relations
in its trace-search input}. The independent trace checker can replay equality
terms if a future producer supplies them, but that larger search was not
implemented. This producer/checker asymmetry is intentional and recorded.

The elementary implementation enumerates rotations and costs quadratic time in
a word's length. More sophisticated minimal-rotation and necklace indexing can
be an optimization later; it would not alter the explicit commutator receipt.

\subsection{An exact separation between operator and trace conclusions}

For self-adjoint $x,y$, let
\[
 h=[x,[x,y]]=x^2y-2xyx+yx^2,\qquad p=I+2h.
\]
The polynomial $h$ is self-adjoint and has zero trace, so $\tr(p)=n\ge0$ for
$n\times n$ matrices. The trace producer certifies it using the single square
$1^*1$ and explicit cyclic residuals.

Take
\[
 x=\begin{pmatrix}1&0\\0&0\end{pmatrix},\qquad
 y=\begin{pmatrix}0&1\\1&0\end{pmatrix}.
\]
Then $h=y$ and
\[
 p=\begin{pmatrix}1&2\\2&1\end{pmatrix},\qquad
 \tr(p)=2,\qquad (1,-1)p(1,-1)^{\mathsf T}=-2.
\]
It is not positive semidefinite. Changing the receipt kind from trace to
operator must not convert this successful trace proof into a matrix inequality.
The strict kind/field checks and an exact quadratic-form test enforce this
boundary in the delivered suite.

\subsection{Why not turn cyclic equivalence into rewriting everywhere?}

Cyclic rotation is legitimate under a trace, but not under arbitrary matrix
contexts. From $\tr(AB)=\tr(BA)$ one cannot conclude
$Q^*ABQ=Q^*BAQ$, or cancel a trace wrapper and return an equality to
\code{grind}. Therefore the trace worker exports only a scalar trace theorem.
A future typed frontend should preserve the observation wrapper as part of the
obligation identity, rather than treating it as irrelevant syntax.
